% Embedded Xinu Raspberry Pi 5 — User's Manual (English)
% Build: pdflatex xinu-pi5-manual-en.tex
\documentclass[a4paper,11pt]{article}

\usepackage[margin=24mm]{geometry}
\usepackage{listings}
\usepackage{xcolor}
\usepackage{enumitem}
\usepackage{longtable}
\usepackage[hidelinks]{hyperref}

\definecolor{codebg}{gray}{0.96}
\lstset{
  basicstyle=\ttfamily\small,
  backgroundcolor=\color{codebg},
  frame=single,
  rulecolor=\color{gray!40},
  breaklines=true,
  columns=fullflexible,
  keepspaces=true,
  showstringspaces=false,
  xleftmargin=4pt, xrightmargin=4pt,
  aboveskip=6pt, belowskip=6pt,
}
\newcommand{\cmd}[1]{\texttt{#1}}

\title{\textbf{Embedded Xinu\\ Raspberry Pi 5 (BCM2712) \\ User's Manual}}
\author{Xinu-rpi5 Project}
\date{June 2026}

\begin{document}
\maketitle
\tableofcontents
\newpage

%======================================================================
\section{Introduction}
This document is the user's manual for the Embedded Xinu kernel
(\texttt{xinu-rpi5}) that runs bare-metal on the Raspberry Pi 5
(BCM2712 / Cortex-A76). The system is a single-image AArch64 kernel with
the MMU enabled, and provides:

\begin{itemize}[noitemsep]
  \item HDMI framebuffer and a window system (desktop)
  \item Wired Ethernet over RP1 (TCP/IP and an HTTP gateway)
  \item Full WiFi via the on-board CYW43455 (scan / WPA2 / DHCP / ping)
  \item USB mouse and keyboard via the RP1 xHCI
  \item microSD read/write (FAT32) mounted at \texttt{/microsd}
  \item A \texttt{kexec} selector that switches to another kernel in RAM
  \item Four OS variants derived from a single source tree
\end{itemize}

\paragraph{Important prerequisite (boot media).}
The board \textbf{boots from a USB stick}; the on-board microSD slot is
normally empty. The firmware loads \texttt{kernel\_2712.img} from the
USB stick's FAT partition (\texttt{bootfs}). The microSD slot is used as
a \textbf{data} area handled by this kernel's SD driver
(mounted at \texttt{/microsd}).

%======================================================================
\section{System Overview}

\subsection{Key Specifications}
\begin{longtable}{ll}
\hline
Item & Value \\
\hline
SoC & BCM2712 (Cortex-A76, AArch64) \\
Kernel image & \texttt{kernel\_2712.img} (entry \texttt{0x80000}) \\
MMU & Enabled (identity map, caches off) \\
Debug UART & \texttt{0x107D001000} (3-pin JST, 115200 8N1) \\
HDMI & Default 1920$\times$1080$\times$32 (OS3 uses 1280$\times$720) \\
Wired IP (static) & \texttt{192.168.3.101} \\
microSD controller & sdhci-brcmstb sdio1 (\texttt{0x10\_00FFF000}) \\
USB host & RP1 DWC3/xHCI (over PCIe) \\
\hline
\end{longtable}

\subsection{Console}
Output appears on both the HDMI text console and the debug UART
(\texttt{uart\_putc} is mirrored to the screen). When the desktop starts,
the window system takes over the screen and the text console is hidden.
Remote operation over the network (Section~\ref{sec:remote}) is the most
reliable way to drive the system.

%======================================================================
\section{Flashing and Booting the Kernel}

\subsection{Build}
Run \texttt{make} in the \texttt{compile/} directory.

\begin{lstlisting}[language=bash]
cd compile
make pi5        # Full OS         -> kernel_2712.img
make pi5-osmin  # OS1 (minimal)   -> kernel_min.img
make pi5-os2    # OS2 (no WM)      -> kernel_os2.img
make pi5-os3    # OS3 (lower res)  -> kernel_os3.img
\end{lstlisting}

\subsection{Writing to the USB Stick (Mac)}
Insert the USB stick (FAT partition = \texttt{bootfs}) into the Mac, copy
the image, and eject. Always verify the mount point.

\begin{lstlisting}[language=bash]
diskutil mount /dev/disk4s1
cp kernel_2712.img /Volumes/bootfs/kernel_2712.img
sync
diskutil unmount /dev/disk4s1
\end{lstlisting}

After writing, put the USB stick back into the Pi 5 and power it on. The
boot log appears on the UART/HDMI after the firmware initialises (about
ten-plus seconds).

%======================================================================
\section{The Four OS Variants and Switching with kexec}
\label{sec:variants}

Four kernels are produced from the same \texttt{main.c} by changing only
the compile flags. Any of them can be switched at runtime via
\texttt{kexec} (Section~\ref{sec:kexec}).

\begin{longtable}{lll}
\hline
Name & Image & Characteristics \\
\hline
Full OS & \texttt{kernel\_2712.img} & Window system + networking (all features) \\
OS1 & \texttt{OS1.IMG} & Shell only. No networking, no windows \\
OS2 & \texttt{OS2.IMG} & All features but no window system (text console) \\
OS3 & \texttt{OS3.IMG} & All features, HDMI lowered one rank to 1280$\times$720 \\
\hline
\end{longtable}

\begin{itemize}[noitemsep]
  \item \textbf{OS1 (minimal)}: no networking and no windows. Shows an
        \texttt{xinu-min\$} prompt on the HDMI text console; driven by the
        USB keyboard.
  \item \textbf{OS2 (no WM)}: keeps all features (networking, etc.) but runs
        the shell on a full-screen text console instead of the window manager
        (prompt \texttt{xinu-os2\$}). Remote operation over the network also
        works.
  \item \textbf{OS3 (lower resolution)}: same as the full OS but with HDMI at
        1280$\times$720.
\end{itemize}

\paragraph{Readability.} The text-console font is the 8$\times$8 glyph scaled
3$\times$ (active on OS1 and OS2).

%======================================================================
\section{microSD Storage}
\label{sec:storage}

The on-board microSD slot is mounted at \texttt{/microsd} (the directory
exists even with no card inserted). Reading a FAT32 card and persistently
writing small files are both supported.

\subsection{Listing and Reading}
\begin{lstlisting}[language=bash]
ls /microsd                 # list files on the card
cat /microsd/CONFIG.TXT     # print a file's contents
\end{lstlisting}

\subsection{Writing (persistent)}
A small file (up to one cluster, 32\,KB) can be written persistently to the
card over HTTP. The body is the content; the end of the URL is the file name
(8.3 form).

\begin{lstlisting}[language=bash]
curl -X POST --data-binary "hello from xinu" \
     http://192.168.3.101/microsd/write/TEST.TXT
# Survives a reboot: ls /microsd / cat /microsd/TEST.TXT
\end{lstlisting}

\subsection{Diagnostics}
A diagnostic command that shows the SD controller and card init state.
\begin{lstlisting}[language=bash]
sdtest        # shows controller + CMD0/CMD8/ACMD41 responses
\end{lstlisting}

%======================================================================
\section{kexec Kernel Selector}
\label{sec:kexec}

\texttt{kexec} loads a kernel image from the microSD into RAM and jumps to
it without a power cycle (no re-flash). Because it is a RAM-only operation,
even a bad image is recovered by power-cycling back into the USB full OS;
there is no brick risk.

\begin{lstlisting}[language=bash]
kexec /microsd/OS1.IMG      # boot OS1 (minimal)
kexec /microsd/OS2.IMG      # boot OS2 (no WM)
kexec /microsd/OS3.IMG      # boot OS3 (lower resolution)
kexec /microsd/KERNEL~1.IMG # the full OS stored on the card
\end{lstlisting}

\paragraph{Notes.}
\begin{itemize}[noitemsep]
  \item Networking drops after \texttt{kexec} (PCIe is re-initialised from a
        non-firmware state). Check the result on the HDMI screen.
  \item Only \textbf{bare-metal Xinu kernels} can be booted. A Linux kernel
        (e.g.\ \texttt{KERNEL8.IMG}, about 9.6\,MB) needs a different boot
        protocol and would hang under the simple chainloader. This
        implementation refuses any image $\geq$ 4\,MB as Linux (both Xinu and
        Linux are arm64 Images sharing the same magic, so size is used to tell
        them apart).
\end{itemize}

%======================================================================
\section{Shell Commands}
Run from the on-screen shell (USB keyboard) or over the network
(Section~\ref{sec:remote}) via \texttt{/run?cmd=...}. The main commands:

\begin{longtable}{ll}
\hline
Command & Description \\
\hline
\cmd{help} & list commands \\
\cmd{ls [path]} & list a directory \\
\cmd{cat <file>} & print a file's contents \\
\cmd{cd / pwd} & change / print the current directory \\
\cmd{kexec /microsd/<k.img>} & boot another kernel (Section~\ref{sec:kexec}) \\
\cmd{cc <file> / make <file>} & compile and run a C/AIPL program in place \\
\cmd{mem} & heap / BSS status \\
\cmd{peek <hex>} & read a 32-bit MMIO word \\
\cmd{uptime} & generic timer value \\
\cmd{ps} & core / EL status \\
\cmd{wifi on/off/status/scan} & connect / disconnect / state / scan (Section~\ref{sec:wifi}) \\
\cmd{wifi adhoc <ssid> [ch] [n]} & join a MANET ad-hoc (mesh) cell (Section~\ref{sec:wifi}) \\
\cmd{wifi aodv <a.b.c.d>} & discover a multi-hop mesh route (Section~\ref{sec:wifi}) \\
\cmd{sdtest} & SD controller diagnostics \\
\cmd{reboot} & reboot (PM watchdog) \\
\cmd{clear} & clear the screen \\
\hline
\end{longtable}

%======================================================================
\section{Remote Operation over the Network}
\label{sec:remote}
The Full OS, OS2, and OS3 listen for HTTP on the wired IP
\texttt{192.168.3.101}. Even where serial input is awkward, the system can
be driven with \texttt{curl} from a Mac or similar.

\begin{longtable}{ll}
\hline
Endpoint & Description \\
\hline
\cmd{/run?cmd=<shell>} & run any shell command and return its output \\
\cmd{/fs}, \cmd{/fs/ls/<d>}, \cmd{/fs/cat/<f>} & browse the VFS tree \\
\cmd{/fs/write/<f>} (POST) & write a file to the VFS (tmpfs) \\
\cmd{/microsd/write/<NAME>} (POST) & persistent write to microSD (Section~\ref{sec:storage}) \\
\cmd{/usb/...} & USB enumeration and diagnostics (Section~\ref{sec:usb}) \\
\cmd{/api/actors}, \cmd{/send?to=\&m=} & actor / AIPL gateway \\
\hline
\end{longtable}

\begin{lstlisting}[language=bash]
curl 'http://192.168.3.101/run?cmd=ls%20/microsd'
curl 'http://192.168.3.101/run?cmd=cat%20/microsd/CONFIG.TXT'
curl 'http://192.168.3.101/run?cmd=wifi%20status'
\end{lstlisting}

\noindent\textbf{Note}: encode spaces in the URL as \texttt{\%20} or
\texttt{+}. The output buffer is bounded, so long results are split.

%======================================================================
\section{WiFi}
\label{sec:wifi}
Full WiFi via the on-board CYW43455 (scan / WPA2 / DHCP / ping) is supported.
The system \textbf{does not auto-connect at boot}; it connects only when
\cmd{wifi on} is run.

\subsection{Connecting to an access point}
\begin{lstlisting}[language=bash]
wifi on <ssid> <pass>   # bring up the firmware (~10 s), join WPA2, run DHCP
wifi on                 # reuse the last creds, or the build-time wifi.conf
wifi scan               # bring up the radio and list nearby APs
wifi status             # connection state, SSID, IP
wifi off                # radio down
wifi-invest             # diagnostics when the link won't come up
\end{lstlisting}
A successful \cmd{wifi on} ends with \texttt{wifi: CONNECTED IP=...}. A WiFi
signal icon is drawn at the bottom-right of the screen, with the SSID and local
IP beneath it. Credentials are embedded into the kernel at build time from a
\texttt{.gitignore}d \texttt{wifi.conf} (never commit the password).

\subsection{Mesh networking (MANET ad-hoc)}
Several Pi 5 (and Pi 4 / Pi 3) Xinu nodes can form a peer-to-peer mesh with
\textbf{no access point}, using 802.11 IBSS (ad-hoc) mode plus on-demand AODV
multi-hop routing.

\begin{lstlisting}[language=bash]
wifi adhoc <cell-ssid> [ch] [node]   # join an ad-hoc cell as 10.0.0.<node>
wifi aodv  <a.b.c.d>                 # discover a multi-hop route on demand
\end{lstlisting}

\begin{itemize}[noitemsep]
  \item \cmd{wifi adhoc mesh1 6 1} brings the radio up in IBSS mode on the named
        cell (BSSID \texttt{02:4d:41:4e:45:54} = ``MANET''), channel 6 (default),
        and takes the static IP \texttt{10.0.0.1} (the \texttt{[node]} number,
        default 1); the AODV relay then runs automatically.
  \item All nodes use the \textbf{same cell SSID and channel, with a distinct
        \texttt{[node]} number} (\texttt{10.0.0.<node>/24}). Nodes in radio range
        reach each other directly at \texttt{10.0.0.x}.
  \item For a destination beyond direct range, \cmd{wifi aodv 10.0.0.3}
        broadcasts an RREQ (UDP 654); intermediate nodes relay it and the
        destination replies with an RREP, installing a multi-hop route
        (e.g.\ \texttt{AODV route: 10.0.0.3 via 10.0.0.2, 2 hop}). Run
        \cmd{wifi adhoc} first so the node has an IP.
\end{itemize}
Ad-hoc/AODV is independent of the infrastructure \cmd{wifi on} mode and is not
restored after reboot --- re-run \cmd{wifi adhoc} on each node.

%======================================================================
\section{USB (Mouse and Keyboard)}
\label{sec:usb}
The USB-A ports connect to two DWC3/xHCI controllers inside the RP1.
To avoid faulting the boot, enumeration runs automatically a few seconds
after boot (or can be driven manually via the \texttt{/usb/...} endpoints).
The boot mouse moves the cursor and the keyboard feeds the shell.

%======================================================================
\section{Appendix}

\subsection{Build Targets}
\begin{longtable}{ll}
\hline
Target & Output \\
\hline
\cmd{make pi5} & \texttt{kernel\_2712.img} (Full OS) \\
\cmd{make pi5-osmin} & \texttt{kernel\_min.img} (OS1) \\
\cmd{make pi5-os2} & \texttt{kernel\_os2.img} (OS2) \\
\cmd{make pi5-os3} & \texttt{kernel\_os3.img} (OS3) \\
\hline
\end{longtable}

\subsection{Troubleshooting}
\begin{itemize}[noitemsep]
  \item \textbf{HTTP not responding}: right after boot the network is not yet
        up. Wait ten-plus seconds and retry. Networking drops after
        \texttt{kexec}; power-cycle to return to the USB full OS.
  \item \textbf{\texttt{kexec} says ``too large''}: an image $\geq$ 4\,MB
        (e.g.\ a Linux kernel) cannot be chainloaded. Specify a bare-metal
        Xinu image.
  \item \textbf{\texttt{/microsd} is empty}: check that a microSD is in the
        slot. Use \texttt{sdtest} to inspect the init state.
  \item \textbf{No display}: if the HDMI framebuffer mailbox init fails, the
        console is UART-only.
\end{itemize}

\subsection{Safety Notes}
\begin{itemize}[noitemsep]
  \item Never commit the WiFi password to the source repository (it is
        embedded at build time from a \texttt{.gitignore}d file).
  \item \texttt{kexec} is a RAM-only operation; even an invalid image is
        recovered by power-cycling (no risk of flash corruption).
\end{itemize}

\vfill
\begin{center}\small Embedded Xinu / Raspberry Pi 5 --- June 2026\end{center}

\end{document}
